\chapter{Descrizione dei file principali}

In questa sezione vengono descritti i file fondamentali del backend Ruggine, con riferimenti puntuali al codice.

\section{main.rs}
File di ingresso dell'applicazione. Si occupa di avviare il server, configurare lo stato globale e inizializzare i servizi principali.

\begin{lstlisting}[language=Rust, caption={Avvio del server in main.rs}]
// ...estratto da src/main.rs...
#[tokio::main]
async fn main() {
    // Inizializzazione logging, stato, router
    // Avvio server Axum
}
\end{lstlisting}

\section{router.rs}
Definisce tutte le rotte HTTP e WebSocket dell'applicazione, collegandole ai rispettivi controller.

\section{state.rs}
Gestisce lo stato condiviso tra le varie componenti (es. pool database, canali di comunicazione).

\section{error.rs}
Definisce le tipologie di errore gestite globalmente dal backend.

\section{controllers/}
Contiene i controller che gestiscono le richieste in ingresso e orchestrano la logica applicativa.

\section{services/}
Contiene la logica di business, separando le regole applicative dall'accesso ai dati.

\section{repositories/}
Gestisce l'accesso al database e l'esecuzione delle query.

\section{DTOs/}
Definisce le strutture dati per la comunicazione tra livelli e verso l'esterno.

\section{migrations/}
Script SQL per la creazione e modifica delle tabelle nel database.